% The page of full-width figures before this one leaves LaTeX with a short height for the
% next column, which would push the first example into the second column. Reset it.
\makeatletter\global\@colroom\@colht\global\vsize\@colht\makeatother
\section{Worked Examples}
\label{app:examples}

% Worked-example boxes. Macro namespace: \wx* only.
\definecolor{wxright}{HTML}{1BAF7A}
\definecolor{wxwrong}{HTML}{D95926}
\definecolor{wxhold}{HTML}{C8891B}
\definecolor{wxrule}{HTML}{BFBFBF}
\definecolor{wxtitle}{HTML}{EFEFEF}
\newcommand{\wxR}{\textcolor{wxright}{$\checkmark$}}
\newcommand{\wxW}{\textcolor{wxwrong}{$\times$}}
\newcommand{\wxA}{\textcolor{wxhold}{$\circ$}}
\newtcolorbox{wxbox}[3]{%
  breakable, enhanced, sharp corners, boxrule=0.4pt,
  colback=white, colframe=wxrule,
  left=4pt, right=4pt, top=3pt, bottom=3pt,
  toptitle=1pt, bottomtitle=1pt,
  colbacktitle=wxtitle, coltitle=black, titlerule=0pt,
  fonttitle=\footnotesize, halign title=flush left,
  title={\textbf{#1}\hspace{0.6em}\textbf{#2}\hfill\texttt{\scriptsize #3}},
}
\newcommand{\wxmeta}[1]{{\scriptsize\textcolor{black!62}{#1}\par}\vspace{1.5pt}}
\newcommand{\wxq}[1]{{\footnotesize\textbf{Q}\hspace{0.5em}#1\par}}
\newcommand{\wxgold}[1]{{\footnotesize\textbf{Gold}\hspace{0.5em}#1\par}}
\newcommand{\wxsep}{\vspace{2.5pt}{\color{wxrule}\hrule height 0.4pt}\vspace{2.5pt}}
\newcommand{\wxrow}[3]{%
  \par\vspace{1.5pt}%
  \noindent%
  \parbox[t]{0.300\linewidth}{\footnotesize\raggedright #1}%
  \hspace{1pt}%
  \parbox[t]{0.045\linewidth}{\footnotesize #2}%
  \hspace{1pt}%
  \parbox[t]{0.620\linewidth}{\footnotesize\raggedright #3}\par}
\newcommand{\wxnote}[1]{{\footnotesize #1\par}}
\newcommand{\wxsnip}[1]{%
  \par\vspace{2pt}\noindent\hspace{0.02\linewidth}%
  \parbox{0.96\linewidth}{\scriptsize\ttfamily\raggedright\sloppy #1}\par\vspace{2pt}}
\newcommand{\wxread}[1]{\wxsep{\footnotesize\textit{#1}\par}}

This appendix shows individual questions behind the rates of \S\ref{sec:findings}, at least one for each failure mechanism of \S\ref{sec:attribution_framework}. Each box gives the question, the gold answer, and the reasoner's output under every condition shown, marked \wxR{} correct, \wxW{} incorrect or \wxA{} abstained. The reasoner is Qwen3-VL-8B at \textbf{TLV} with no instruction unless a row says otherwise, and verdicts follow the accuracy rubric. Page numbers are zero-based indices into the document. None of these questions is one of the six in Figure~\ref{fig:mechanism_examples}.

%% ---------------------------------------------------------------------------
\subsection{Representation}
\label{app:ex_representation}

\begin{wxbox}{R1}{Modality ceiling}{mmlongbench:000014}
\wxmeta{Presentation deck (Tutorial/Workshop), born-digital. Gold page 13, a single chart. Evidence source: chart.}
\wxq{According to the chart on page 14 how much time was spent with family and friends in 2010?}
\wxgold{21\%}
\wxsep
\wxnote{The page is one chart, and the parser has nothing to transcribe. PaddleOCR-VL emits only a placeholder:}
\wxsnip{<div style="text-align: center;"><img src="[\ldots]" alt="Image" width="52\%" /></div>}
\wxsep
\wxrow{T}{\wxA}{``The document does not provide any evidence or chart on page 14, so it is not possible to determine how much time was spent with family and friends in 2010 [\ldots]''}
\wxrow{TL}{\wxA}{``[\ldots] The image referenced (on page 14) is not described or transcribed, so its content cannot be interpreted. [\ldots]''}
\wxrow{V}{\wxR}{``According to the chart on page 14, in 2010, 21\% of weekend time was spent with family and friends.''}
\wxrow{TLV}{\wxR}{``According to the chart on page 14, in 2010, 21\% of weekend time was spent with family and friends.''}
\wxread{The gold page reaches every representation, so nothing is missing from the context. The value lives in the chart's geometry and no text channel carries it, so both text representations decline instead of guessing, which is the refusal pattern of Table~\ref{tab:abstention_rungs}.}
\end{wxbox}

\begin{wxbox}{R2}{Conversion fidelity: parser}{mmlongbench:000513}
\wxmeta{Guidebook, born-digital. Gold page 9, a numbered list of market names with a code beside each. Evidence source: text and table.}
\wxq{How many geographic market names are listed under Los Angeles area?}
\wxgold{13}
\wxsep
\wxnote{The list runs from item 14 to item 26. The three parsers keep its text and differ in what they do with the code column:}
\wxsnip{PaddleOCR-VL: California 14-26 \textbar{} 14. San Fernando Valley (West) CA14 \textbar{} [\ldots] \textbar{} 26. Santa Ana CA26 \\[1.5pt] MinerU 2.5: 14. San Fernando Valley (West) \textbar{} [\ldots] \textbar{} 26. Santa AnaEPS \textbar{} Code \textbar{} 4.14 \textbar{} [\ldots] \textbar{} 4.26 \\[1.5pt] Unlimited-OCR: <table><tr><td>14.</td><td>San Fernando Valley (West)</td><td>CA14</td></tr> [\ldots]}
\wxsep
\wxrow{T}{\wxR}{``The document lists \textbf{13} geographic market names under the Los Angeles area. [\ldots]''}
\wxrow{TL, PaddleOCR-VL}{\wxW}{``26 The document lists 26 geographic market names under the `Los Angeles Area' section, ranging from CA14 to CA26.''}
\wxrow{TL, MinerU 2.5}{\wxW}{``26 The document lists 26 geographic market names [\ldots] ranging from `San Fernando Valley (West)' (item 14) to `Santa Ana' (item 26).''}
\wxrow{TL, Unlimited-OCR}{\wxR}{``There are \textbf{13} geographic market names listed under the Los Angeles area. The list includes entries numbered 14 through 26 [\ldots]''}
\wxrow{TLV, PaddleOCR-VL}{\wxR}{``[\ldots] there are \textbf{13} geographic market names listed under the Los Angeles area. [\ldots] codes from 14 to 26. Counting these entries [\ldots]''}
\wxread{Every parser returns all thirteen names, yet two flatten the list so that its last label reads as its length. The parser that keeps the rows as a table gives the right count, as does the same flattened parse once the page image accompanies it: the image is the redundant channel that absorbs the conversion error.}
\end{wxbox}

\begin{wxbox}{R3}{Conversion fidelity: resolution}{mmlongbench:000065}
\wxmeta{Academic paper, born-digital. Gold page 6, a results table with one row per model and temperature. Evidence source: text and table.}
\wxq{From the paper, which temperature gives ChatGPT the highest alignment score?}
\wxgold{0.1}
\wxsep
\wxrow{V, \texttt{low}}{\wxW}{``[\ldots] the temperature that gives ChatGPT the highest alignment score is \textbf{0.5}.''}
\wxrow{V, \texttt{med}}{\wxW}{``[\ldots] is \textbf{0.5}. [\ldots] Table 3 lists three ChatGPT models: ChatGPT (0.5), ChatGPT (0.1), and ChatGPT (0.9) [\ldots]''}
\wxrow{V, \texttt{high}}{\wxR}{``[\ldots] the temperature that gives ChatGPT the highest alignment score is \textbf{0.1}.''}
\wxrow{TLV, \texttt{low}}{\wxR}{``[\ldots] ChatGPT (0.1): 85.9, ChatGPT (0.5): 84.5, ChatGPT (0.9): 84.1. [\ldots] \textbf{85.9} is the highest alignment score, which corresponds to [\ldots] \textbf{0.1}.''}
\wxread{The image alone finds the right table at every resolution but only reads its small digits correctly at \texttt{high}. With parser text the values are available as characters, so the lowest resolution already suffices, as in the table rows of Table~\ref{tab:app_resolution_sources}.}
\end{wxbox}

%% ---------------------------------------------------------------------------
\subsection{Selection}
\label{app:ex_selection}

\begin{wxbox}{S1}{Evidence coverage}{mmlongbench:000082}
\wxmeta{Survey paper (Academic paper), born-digital. Gold pages 8 and 9: the sentence that answers the question starts on the last line of page 8 and finishes on page 9. \texttt{drop\_top} leaves page 8 and \texttt{drop\_bottom} page 9. Evidence source: text.}
\wxq{Which method integrates knowledge into white-box models via directive fine-tuning?}
\wxgold{PKG}
\wxsep
\wxrow{gold pages}{\wxR}{``[\ldots] the method that integrates knowledge into white-box models via directive fine-tuning is \textbf{PKG}. [\ldots] `PKG introduces an innovative method for integrating knowledge into white-box models via directive fine-tuning [75].' [\ldots]''}
\wxrow{\texttt{drop\_top}}{\wxW}{``[\ldots] the method that integrates knowledge into white-box models via directive fine-tuning is \textbf{Fine-tuning Embedding Model}. [\ldots]''}
\wxrow{\texttt{drop\_bottom}}{\wxW}{``The method that integrates knowledge into white-box models via directive fine-tuning is \textbf{directive fine-tuning}. The document explicitly states: `introduces an innovative method for integrating knowledge into white-box models via directive fine-tuning [75].' [\ldots]''}
\wxread{Page 8 ends on the method's name and page 9 opens with what it does, so neither page holds the answer. With one of them withheld the reasoner does not decline: it names a nearby section heading, or repeats the question's own phrase. Both removals break the answer, so retrieval rank does not identify a page that can safely be omitted.}
\end{wxbox}

\begin{wxbox}{S2}{Distractor exposure: broken}{mmlongbench:000078}
\wxmeta{Survey paper (Academic paper), born-digital, multi-hop. Gold pages 5 (a table of methods) and 16 (the reference list). $+3$ adds pages 0, 6 and 7; page 0 is the survey's own title page. 4-bit weights.}
\wxq{What is the paper's full title that proposes the method that has a retrieval granularity of phrase?}
\wxgold{Copy is all you need.}
\wxsep
\wxrow{gold only}{\wxR}{``[\ldots] the method with a retrieval granularity of ``Phrase'' is \textbf{CoG [29]}. [\ldots] the full title of the paper proposing the CoG method is: \textbf{``Copy is all you need''}. This is found in the References section [\ldots]''}
\wxrow{$+3$}{\wxW}{``[\ldots] the paper's full title is \textbf{Retrieval-Augmented Generation for Large Language Models: A Survey}. This is explicitly stated in the title section of the document [\ldots]''}
\wxrow{$+5$ to $+18$}{\wxW}{The same answer at every completed depth.}
\wxread{Both gold pages are still present and nothing in them changed. The added title page is a plausible partial match for ``the paper's full title'', and the reasoner takes it instead of following the table entry to its reference. Coverage is complete; the failure is in using the evidence amid competing context.}
\end{wxbox}

\begin{wxbox}{S3}{Distractor exposure: repaired}{mmlongbench:000243}
\wxmeta{Rezoning study (Research report), born-digital, single-hop. Gold page 7, which carries a table of residential capacity by borough. $+3$ adds pages 4, 8 and 12, the surrounding discussion. 4-bit weights.}
\wxq{What is the residential capacity of Staten Island from 2003 to 2007? Give me an integer.}
\wxgold{435000000}
\wxsep
\wxrow{gold only}{\wxW}{``80000''}
\wxrow{$+3$}{\wxR}{``435000000''}
\wxrow{$+10$}{\wxR}{``435000000''}
\wxrow{$+13$}{\wxW}{``5980000''}
\wxrow{$+18$}{\wxR}{``435000000''}
\wxread{The table holding the answer is supplied in every row. The reasoner misreads it in isolation and reads it correctly once neighbouring pages of the same section are present, then fails again at one deeper setting. Added pages therefore both repair and break answers on the same question, the two-way movement behind the transitions of Figure~\ref{fig:distractor_k}.}
\end{wxbox}

\paragraph{Distractors on LongDocURL.}
Table~\ref{tab:app_ldu_distractor} repeats the distractor sweep on LongDocURL with Qwen3-VL-8B. This experiment is still running, and the table reports the depths completed so far. On single-hop questions the first three added pages lower accuracy from 77.2\% to 68.5\%, breaking 13.5\% of answers and repairing 4.8\%, close to the first step on MMLongBench-Doc.

%% ---------------------------------------------------------------------------
\subsection{Reasoning}
\label{app:ex_reasoning}

\begin{wxbox}{N1}{Evidence integration}{mmlongbench:001015}
\wxmeta{Annual report (Financial report), born-digital, multi-hop. Gold pages 35 (the income statement) and 39 (the cash flow statement). Rows are Qwen3-VL sizes. Evidence source: table.}
\wxq{what is EBITDA for costco in FY2021?}
\wxgold{8489}
\wxsep
\wxrow{2B}{\wxW}{``[\ldots] EBITDA for Costco in FY2021 is \$6,708 million. This is found in the [\ldots] row labeled `Operating income' [\ldots]''}
\wxrow{4B}{\wxW}{``The EBITDA for Costco in FY2021 is \$6,708 million. [\ldots] EBITDA [\ldots] is typically calculated as Operating Income plus Depreciation and Amortization, but in this document, the `Operating income' row is presented as the direct measure of EBITDA [\ldots]''}
\wxrow{8B}{\wxW}{``[\ldots] Operating Income + Depreciation and Amortization + Interest Expense + Provision for Income Taxes = \$6,708 + \$1,781 + \$(171) + \$1,601 = \textbf{\$8,469 million}''}
\wxrow{32B}{\wxR}{``[\ldots] From the Consolidated Statements of Income: Operating Income = \$6,708 million. From the Consolidated Statements of Cash Flows: Depreciation and Amortization = \$1,781 million. [\ldots] \$6,708 + \$1,781 = \textbf{\$8,489 million}''}
\wxread{Every model receives the same two statements. The two smallest stop at operating income on the first page, one of them after stating the formula that needs the second. The 8B reaches the cash flow statement but adds lines that operating income already excludes, and only the largest joins the two operands. The evidence was sufficient throughout; what scales is the ability to combine it.}
\end{wxbox}

\begin{wxbox}{N2}{Response calibration: answerable}{mmlongbench:000418}
\wxmeta{E-commerce market deck (Research report), scanned, multi-hop. Gold pages 11 and 14, one bar chart on each. Evidence source: chart.}
\wxq{How many people in India were using a debit card when the e-commerce sales reached 11 billion?}
\wxgold{399 million}
\wxsep
\wxrow{\texttt{none}}{\wxR}{``[\ldots] when e-commerce sales reached \$11 billion (in 2014), the number of debit card users in India was \textbf{399 million}. [\ldots]''}
\wxrow{\texttt{cot}}{\wxR}{``Step 1: [\ldots] the \$11 billion mark is shown for the year 2014. Step 2: [\ldots] the bar chart titled `Number of Debit Card users in India (In millions)' shows the value for 2014 as 399 million. [\ldots] Answer: 399 million''}
\wxrow{\texttt{abstain}}{\wxA}{``Not answerable''}
\wxrow{\texttt{abstain\_cot}}{\wxR}{``Step 1: [\ldots] e-commerce sales were \$11 billion in 2014. Step 2: [\ldots] For 2014, the number is 399 million. [\ldots] Answer: 399 million''}
\wxread{The answer is on neither page alone: one chart gives the year, the other the count for that year. Under \texttt{abstain} the reasoner refuses a question it answers correctly with no instruction, a false refusal. Asking it to reason before deciding makes it locate both values first, and the refusal disappears.}
\end{wxbox}

\begin{wxbox}{N3}{Response calibration: unanswerable}{mmlongbench:000047}
\wxmeta{Data-leadership talk (Tutorial/Workshop), scanned, unanswerable. BM25's top three pages (0, 1, 2) are supplied; the document mentions no video.}
\wxq{Which Youtube does the slides use to show the consequce of having a weak data leadership?}
\wxgold{Not answerable}
\wxsep
\wxrow{\texttt{none}}{\wxW}{``The slides use a YouTube clip from the TV show \textit{Breaking Bad} to illustrate the consequence of having weak data leadership [\ldots] as symbolized by the character Walter White hoarding money [\ldots]''}
\wxrow{\texttt{cot}}{\wxR}{``[\ldots] Step 4: However, the document evidence does not mention any YouTube or reference any video content. The image is a still image, not a YouTube video. [\ldots] Answer: Not specified in the document.''}
\wxrow{\texttt{abstain}}{\wxR}{``Not answerable''}
\wxrow{\texttt{abstain\_cot}}{\wxR}{``[\ldots] Step 3: The document does not mention YouTube at all. [\ldots] Answer: Not answerable''}
\wxread{One supplied page shows a man holding stacks of money, and with no instruction the reasoner builds a show, a character and an interpretation out of it. Each of the other prompts leads it to check the premise against the pages and decline, which is the correct response here.}
\end{wxbox}
z